%% ---------------------------------------------------------------------------
\section{Related work}
\label{sec:related}

\subsection{Certificate size in post-quantum TLS}
\label{sec:related-tls}
That certificates rather than signatures dominate a post-quantum exchange is
established in the TLS literature, and this paper does not claim it as new.
Sikeridis et al.\ found the certificate chain, not the handshake signature,
driving the added latency of post-quantum TLS~1.3, to the point of proposing
different algorithms at different levels of one
chain~\cite{sikeridis2020ndss,sikeridis2020conext}; later studies reached the
same conclusion on newer stacks~\cite{sosnowski2023}, and Cloudflare made the
sizing argument to practitioners~\cite{westerbaan2021}. The remedies all stop
sending certificates the verifier could already have: cached
information~\cite{rfc7924}, abridged certificate compression, which ships a
dictionary of the WebPKI's roots and intermediates so that only the leaf
travels and reports about 75\% off the chain against about 25\% for generic
compression~\cite{certabridge}, and Merkle Tree
Certificates~\cite{merkletreecerts}. Leaf-only carriage is the manifest
equivalent of abridged compression in its simplest form, with the C2PA trust
list as the dictionary, and we present it as such.

\paragraph{What transfers and what does not}
The structural result transfers, and our measurements extend it to Falcon,
SLH-DSA and composites: because every certificate contains a signature, chain
dominance is a property of chains rather than of any scheme. Generic
compression fails here as it does in TLS (Table~\ref{tab:compress}). Four
differences change the conclusions. A manifest chain is paid once per object
and kept for the life of the file, so the cost is storage and bandwidth rather
than handshake latency. Nothing is negotiated, so the carriage form is a
policy for every verifier. A manifest verifier may be offline for good, which
is why we recommend the leaf over the thumbprint where the TLS analogy would
favour the smaller option. And C2PA's trust list lists issuing CAs as anchors,
which makes the leaf-only form conformant without any new mechanism.

\subsection{Post-quantum benchmarking}
\label{sec:related-bench}

Primitive-level performance is well covered, from pqm4 on
microcontrollers~\cite{kannwischer2019pqm4,pqm4,greconici2021,li2024m0} and
NEON and AVX2 implementations on application
processors~\cite{becker2022neon} to batched GPU verification reported in the
millions per second~\cite{shen2024gpu,lee2024falcon,nvidia2024cupqc}. At the
protocol level, Paquin et al.\ benchmarked post-quantum TLS under emulated
loss~\cite{paquin2020}, and the Open Quantum Safe project made such
experiments routine~\cite{stebila2016}. Deployment has followed for key
agreement and not for authentication: by late 2025 most human-initiated
traffic to Cloudflare used post-quantum key agreement while no publicly
trusted post-quantum certificate was in use~\cite{westerbaan2025pq}. The half
of the migration this paper prices has not begun even on the web.

Methodology has become a subject of its own. Riou et al.\ catalogue pitfalls
in ML-DSA benchmarking, among them the tail that rejection sampling gives
signing time~\cite{riou2026}; SUPERCOP has long recorded the compiler and
options with every measurement~\cite{ebacs}, and Mytkowicz et al.\ showed that
measurement setup alone can reverse a conclusion about \texttt{-O2} against
\texttt{-O3}~\cite{mytkowicz2009}. Our audit (Section~\ref{sec:audit}) finds
that applied post-quantum work rarely carries those habits over.
Implementation-controlled measurement is rarer still; our earlier study of HQC
as a fallback to ML-KEM carries such a control for key
establishment~\cite{alfarajat2026hqc}. It shares this paper's method and not
its data or subject: its hardware work runs key establishment over a Thread
link, and the microcontroller work here is signatures only, over Bluetooth LE,
on separate firmware.

\subsection{Content provenance}
\label{sec:related-prov}

C2PA manifests are JUMBF boxes~\cite{iso19566-5} carrying COSE signatures, and
they have left the laboratory: Leica shipped the first camera that signs at
capture in 2023~\cite{cai2023leica}, Sony added signing to its bodies the
following year~\cite{sony2024c2pa}, and Samsung and Google brought content
credentials to phones in 2025~\cite{samsung2025s25,google2025pixel}. Google's
design keeps signing keys in a dedicated security chip, runs a time-stamping
authority on the device and issues a certificate for every
image~\cite{google2025pixel}, the lifecycle under which Section~\ref{sec:cache}
finds caching worth least.

Golaszewski et al.\ give the first independent security analysis
of C2PA and conclude that the specification does not meet its stated
goals~\cite{golaszewski2026eprint,golaszewski2026}. Two of their findings
concern machinery whose post-quantum cost we measure: revocation checking is
optional, and generators and validators disagree about which timestamp is
trusted. The two departures we found in the reference implementation
(Section~\ref{sec:patch}) are of the same kind, in the key-usage rules for
signers and time-stamping authorities. VerITAS proves in zero knowledge that a
published photograph differs from a signed original only by permitted
edits~\cite{datta2025veritas}, and Nemecek et al.\ build images whose valid
manifests contradict their watermarks~\cite{nemecek2026integrity}. Watermarks
and perceptual hashes are the alternatives to a signature; they rest on no
assumption Shor's algorithm breaks, and they fail by removal or
steering~\cite{wen2023treering,gowal2025synthid,an2024waves,zhao2024invisible,%
meta_pdq,photodna,struppek2022,jain2022}.

Post-quantum signatures have been proposed for image authenticity outside this
ecosystem, with zero-knowledge proofs of permitted
transformations~\cite{iavich2025pqveritas} and inside a watermarking framework~\cite{alkhatib2025deepfake}. Neither measures the credential, and to our knowledge this paper is the first published measurement
of post-quantum signatures inside the credential stack that cameras and phones
ship.

\subsection{Migration measurement}
\label{sec:related-migration}

Mixed chains have been measured in TLS: Paul et al.\ placed different
signature algorithms at different levels of a TLS~1.3 chain~\cite{paul2022},
and Bindel et al.\ tested whether hybrid certificates survive real X.509, TLS
and S/MIME implementations~\cite{bindel2017}. Section~\ref{sec:mixedchain}
asks what each ordering protects and what it costs to carry in every object.
Planning for migration is mostly guidance: research agendas~\cite{ott2019},
crypto-agility~\cite{nist_cswp39}, inventory~\cite{nist_sp1800_38,ecma424} and
maturity models~\cite{hohm2023}. Past migrations were slow at the tail: four
years after the MD5 rogue CA~\cite{stevens2009} a trusted CA still issued
MD5-signed certificates~\cite{durumeric2013}, and TLS clients abandon old
algorithms slowly~\cite{kotzias2018}. Provenance has a worse tail than TLS: a
handshake renegotiates on the next connection, and a signed image keeps its
signature until someone signs it again.
